\section*{Hoofdstuk \ref{ch:b2:crustal-melting} --- Smelten van de korst: migmatieten en granieten}

\begin{solution}{exo:b2:crustal-melting:1}
(a) Per \qty{100}{\gram}: $n(\ce{Al2O3}) = 17.03/101.96 = \qty{0.16703}{mol}$,
$n(\ce{CaO}) = 0.09183$, $n(\ce{Na2O}) = 0.06783$, $n(\ce{K2O}) = 0.03076$.
A/CNK $= 0.16703/0.19042 = 0.877$; A/NK $= 0.16703/0.09859 = 1.694$.
(b) A/CNK $< 1 <$ A/NK: \omterm{def:b2:crustal-melting:asi}{metaluminisch}. Het calcium dat de veldspaten niet kunnen opnemen, gaat naar
calciumsilicaten: hoornblende (of clinopyroxeen, titaniet).
\end{solution}

\begin{solution}{exo:b2:crustal-melting:2}
(a) Muscoviet en granaat bevatten aluminium buiten de veldspaten: \omterm{def:b2:crustal-melting:asi}{peraluminisch}; dit graniet
kan van S-type zijn. (b) Hoornblende en titaniet bevatten calcium buiten de veldspaten:
\omterm{def:b2:crustal-melting:asi}{metaluminisch}. (c) Aegirien en een natriumamfibool bevatten natrium zonder het aluminium
van veldspaat: \omterm{def:b2:crustal-melting:asi}{peralkalisch}.
\end{solution}

\begin{solution}{exo:b2:crustal-melting:3}
Een kilometer gesteente weegt $2800\times9.81\times1000 = \qty{2.747e7}{Pa}$, dus
$P = 0.2747\,z$ (\unit{kbar}, $z$ in \unit{km}). De solidus is de lijn $T = 680 - 7.5(P -
2) = 695 - 2.060\,z$. Uit $10 + 20z = 695 - 2.060z$ volgt $z = 685/22.060 =
\qty{31.1}{km}$, $T = \qty{631}{\celsius}$, $P = \qty{8.5}{kbar}$.
\end{solution}

\begin{solution}{exo:b2:crustal-melting:4}
(a) De witte lagen zijn \omterm{def:b2:crustal-melting:migmatite}{leucosomen}, de zwarte naden \omterm{def:b2:crustal-melting:migmatite}{melanosomen}, de gewone gneis het
mesosoom. (b) Afzonderlijke, scherp begrensde \omterm{def:b2:crustal-melting:migmatite}{leucosomen}, \omterm{def:b2:crustal-melting:migmatite}{melanosomen} en mesosomen: een \omterm{def:b2:crustal-melting:migmatite}{metatexiet},
met een stromatische structuur. (c) Het gesteente is een \omterm{def:b2:crustal-melting:migmatite}{diatexiet} geworden (schlierig en nebulitisch):
het smelten ging verder, de donkere mineralen smolten ook, en de smelt bleef ter plaatse en
wiste de gelaagdheid uit.
\end{solution}

\begin{solution}{exo:b2:crustal-melting:5}
(a) $M = 39.10 + 3\times24.31 + 26.98 + 3\times28.09 + 12\times16.00 + 2\times1.008 =
\qty{417.30}{\gram\per\mole}$; $w = 18.02/417.30 = \qty{4.32}{\percent}$.
(b) $F = 0.25\times0.0432/0.040 = 0.270$, of \qty{27}{\percent} smelt in massa: binnen de
20 tot \qty{30}{\percent} in volume die experimenten met \omterm{def:b2:crustal-melting:reactions}{dehydratatiesmelten} van biotiet aan
vergelijkbare gesteenten opleveren.
\end{solution}

\begin{solution}{exo:b2:crustal-melting:6}
(a) Kaliveldspaat \qty{278.35}{\gram\per\mole}, albiet 262.24, anorthiet 278.22, muscoviet
398.33. (b) Hoeveelheden: kaliveldspaat $30/278.35 = \qty{0.10778}{mol}$, albiet $0.13347$,
anorthiet $0.01797$, muscoviet $0.01255$. Kaliveldspaat is $\tfrac12$\ce{K2O.}$\tfrac12$\ce{Al2O3},
albiet $\tfrac12$\ce{Na2O.}$\tfrac12$\ce{Al2O3}, anorthiet \ce{CaO.Al2O3}, muscoviet
$\tfrac12$\ce{K2O.}$\tfrac32$\ce{Al2O3} per formule-eenheid. $n(\ce{Al2O3}) = 0.05389 +
0.06673 + 0.01797 + 0.01883 = \qty{0.15742}{mol}$; $n(\ce{K2O}) = 0.05389 + 0.00628 =
0.06017$, $n(\ce{Na2O}) = 0.06673$, $n(\ce{CaO}) = 0.01797$, som $0.14487$.
A/CNK $= 0.15742/0.14487 = 1.087$. (c) Zonder de muscoviet geven de drie veldspaten
precies 1 (\cref{prop:b2:crustal-melting:feldspar}); de muscoviet, met per formule-eenheid één \ce{Al2O3}
meer dan kaliveldspaat, maakt het gesteente \omterm{def:b2:crustal-melting:asi}{peraluminisch}.
\end{solution}

\begin{solution}{exo:b2:crustal-melting:7}
Basalt: $900 + 2.5\times(1200 - 850) = \qty{1775}{\mega\joule\per\cubic\metre}$. Korst:
$2.5\times(850 - 650) + 0.20\times900 = \qty{680}{\mega\joule\per\cubic\metre}$.
$V_c/V_b = 1775/680 = 2.6$. De schatting verwaarloost de warmte die door geleiding naar boven
en opzij weglekt en de tijd die de warmte nodig heeft om zich te verspreiden; ze neemt ook aan dat alle
warmte van de basalt terechtkomt in korst die klaar is om te smelten: het is een bovengrens.
\end{solution}

\begin{solution}{exo:b2:crustal-melting:8}
(a) $n(\ce{Al2O3}) = 14.15/101.96 = 0.13878$, $n(\ce{CaO}) = 0.01753$, $n(\ce{Na2O}) =
0.07809$, $n(\ce{K2O}) = 0.04416$; A/CNK $= 0.13878/0.13978 = 0.993$, A/NK $=
0.13878/0.12225 = 1.135$: net \omterm{def:b2:crustal-melting:asi}{metaluminisch}. (b) Met \qty{14.30}{\percent} \ce{Al2O3},
$0.14025/0.13978 = 1.003$; met \qty{14.00}{\percent}, $0.13731/0.13978 = 0.982$. De index
ligt binnen de gecertificeerde onzekerheid aan weerszijden van 1: het obsidiaan ligt op de
grens tussen \omterm{def:b2:crustal-melting:asi}{metaluminisch} en \omterm{def:b2:crustal-melting:asi}{peraluminisch}, en zijn klasse is niet robuust.
\end{solution}

\begin{solution}{exo:b2:crustal-melting:9}
(a) Gemiddelden $\bar P = \qty{5.0}{kbar}$, $\bar w = \qty{10.045}{\percent}$;
$\sum(P - \bar P)(w - \bar w) = 17.62$, $\sum(P - \bar P)^2 = 14.0$; helling
$17.62/14.0 = \qty{1.259}{\percent}$ per kilobar, snijpunt $10.045 - 1.259\times5.0 =
\qty{3.75}{\percent}$. De helling stemt tot op \qty{1}{\percent} overeen met de opgegeven
1.27. (b) $3.75 + 1.259\times6 = \qty{11.30}{\percent}$. (c) $10/0.1130 =
\qty{88.5}{\gram}$ smelt. (d) De lijn steunt op gegevens tussen 3 en \qty{8}{kbar} bij alleen
\qty{800}{\celsius}; niets garandeert dat de oplosbaarheid lineair blijft bij hogere
druk, waar de smelt en het fluïdum meer op elkaar gaan lijken.
\end{solution}

\begin{solution}{exo:b2:crustal-melting:10}
(a) De lijn stijgt \qty{4}{kbar} over \qty{50}{\celsius}: $P = 6 + 0.08\,(T - 750)$;
bij \qty{770}{\celsius} is $P = \qty{7.6}{kbar}$. (b) Van 7.6 tot \qty{6}{kbar} daalt de
temperatuur langs de lijn van 770 tot \qty{750}{\celsius}; de vrijgekomen warmte,
$2.5\times20 = \qty{50}{\mega\joule\per\cubic\metre}$, smelt $50/900 = 0.056$, ongeveer
\qty{6}{\percent} van het gesteente in volume. (c) Dit is het venster dat voor de
\omterm{def:b2:crustal-melting:leucogranite}{leucogranieten} van de Himalaya is voorgesteld, 6 tot \qty{8}{kbar} en 750 tot \qty{770}{\celsius}: decompressie
alleen maakt smelt, maar slechts enkele procenten per keer dat de kromme wordt gekruist.
\end{solution}

\begin{solution}{exo:b2:crustal-melting:11}
$\ell^2/\kappa = (2000)^2/10^{-6} = \qty{4e12}{s} = \qty{0.13}{Myr}$ voor de \omterm{def:b2:magma-emplacement:bodies}{sill};
$(10^4)^2/10^{-6} = \qty{1e14}{s} = \qty{3.2}{Myr}$ voor de laag van onderplating. Beide zijn kort
naast tien miljoen jaar: één gebeurtenis kan de korst niet zo lang gesmolten houden; daarvoor is
steeds opnieuw aangevoerde warmte nodig, door herhaalde intrusies of door het trage opwarmen van een
verdikte korst.
\end{solution}

\begin{solution}{exo:b2:crustal-melting:12}
(a) Uit $5.0x + 2.0(1 - x) = 2.75$ volgt $x = 0.75/3.0 = 0.25$: een kwart van het neosoom is
\omterm{def:b2:crustal-melting:migmatite}{leucosoom}. (b) $0.25\times0.20 + 0.75\times4.0 = 0.05 + 3.00 = \qty{3.05}{\percent}$, zoals
in het mesosoom: de interpretatie als gesloten systeem houdt stand. (c) Voor \qty{100}{\gram}
moedergesteente ontstonden \qty{25}{\gram} smelt en \qty{75}{\gram} residu; als $m$ gram
smelt vertrok, geldt $(25 - m)/(100 - m) = 0.10$, dus $25 - m = 10 - 0.1m$ en $m = 15/0.9 =
\qty{16.7}{\gram}$: twee derde van de smelt verliet de ontsluiting.
\end{solution}

\begin{solution}{pb:b2:crustal-melting:1}
\noindent\textbf{Deel I.}

\textbf{1.} $39.10 + 3\times26.98 + 3\times28.09 + 12\times16.00 + 2\times1.008 =
\qty{398.33}{\gram\per\mole}$.

\textbf{2.} $18.02/398.33 = \qty{4.52}{\percent}$.

\textbf{3.} $200\times0.0452 = \qty{9.04}{\gram}$ water per kilogram schist.

\textbf{4.} $9.97 + (13.82 - 9.97)\times2/3 = \qty{12.54}{\percent}$. De kleinste hoeveelheid smelt is
$9.04/0.1254 = \qty{72.1}{\gram}$ per kilogram: $F \ge \qty{7.21}{\percent}$
(\cref{prop:b2:crustal-melting:water}).

\textbf{5.} $9.04/0.060 = \qty{151}{\gram}$: $F = \qty{15.1}{\percent}$.

\textbf{6.} Bijvoorbeeld: het gesteente kan minder muscoviet bevatten; een deel van de muscoviet kan
al onder de solidus, met kwarts, uiteengevallen zijn tot kaliveldspaat en
aluminiumsilicaat, waarbij zijn water als fluïdum vertrok voordat het smelten begon; de smelt kan meer
water bevatten dan \qty{6}{\percent}; wat water kan in biotiet achterblijven.

\noindent\textbf{Deel II.}

\textbf{7.} $680 - 7.5\times(7.8 - 2) = \qty{636.5}{\celsius}$; het gesteente piekt
\qty{43.5}{\celsius} daaronder en kan niet smelten, zelfs niet met vrij water.

\textbf{8.} $605 + (19.2 - 15.3)/(27 - 15.3)\times65 = \qty{626.7}{\celsius}$. Bij
$t = 0$ ligt de basis, op \qty{595}{\celsius}, \qty{32}{\celsius} eronder; bij
\qty{10}{Myr}, op \qty{656}{\celsius}, ligt ze erboven.

\textbf{9.} Smelten bij de natte solidus vraagt vrij water, en een verhit metamorf gesteente
bevat er heel weinig: het vrije water wordt opgebruikt door de eerste smelt, die dan niet meer groeit
tot de mica's uiteenvallen.

\textbf{10.} Tussen 30 en \qty{40}{Myr}: $30 + (800 - 795)/(845 - 795)\times10 =
\qty{31.0}{Myr}$.

\textbf{11.} De basis zakt onder \qty{12.5}{kbar} tussen 70 en \qty{80}{Myr}:
$70 + (13.7 - 12.5)/(13.7 - 12.4)\times10 = \qty{79.2}{Myr}$, bij
$914 - 14\times0.923 = \qty{901}{\celsius}$, binnen het amfiboolinterval: een amfiboliet
aan de basis zou aan het smelten zijn.

\textbf{12.} $90 + (11.0 - 10)/(11.0 - 9.6)\times10 = \qty{97.1}{Myr}$, bij $867 -
53\times0.714 = \qty{829}{\celsius}$. Dat ligt boven het muscovietinterval: de muscoviet
is uiteengevallen. Het ligt binnen de 800 tot \qty{850}{\celsius} van het uiteenvallen van biotiet
(gemeten bij lagere druk): biotiet, het andere waterhoudende mineraal van de schist, zit op zijn
grens.

\noindent\textbf{Deel III.}

\textbf{13.} $n(\ce{Al2O3}) = 15.0/101.96 = \qty{0.14712}{mol}$, $n(\ce{CaO}) =
0.55/56.08 = 0.00981$, $n(\ce{Na2O}) = 3.4/61.98 = 0.05486$, $n(\ce{K2O}) = 4.9/94.20 =
0.05202$.

\textbf{14.} A/CNK $= 0.14712/0.11669 = 1.261$.

\textbf{15.} A/NK $= 0.14712/0.10688 = 1.377$; A/CNK $> 1$: \omterm{def:b2:crustal-melting:asi}{peraluminisch}.

\textbf{16.} \omterm{def:b2:crustal-melting:asi}{Peraluminisch}, uit een schist: een \omterm{def:b2:crustal-melting:types}{graniet van S-type}. Zijn overtollige aluminium zit in
muscoviet en granaat (of cordieriet).

\textbf{17.} $0.14712 - 0.11669 = \qty{0.03043}{mol}$ \ce{Al2O3}; als muscoviet,
$0.03043\times398.33 = \qty{12.1}{\gram}$ in \qty{100}{\gram} graniet.

\textbf{18.} Het obsidiaan is (net) \omterm{def:b2:crustal-melting:asi}{metaluminisch}, zonder het overtollige aluminium van een
metasedimentaire smelt: zijn bron is waarschijnlijk anders. Maar één monster bewijst weinig:
elke granitoïdreeks heeft zwak peraluminische leden, en een reeks wordt beoordeeld op haar hele
trend.

\noindent\textbf{Deel IV.}

\textbf{19.} $10^{12}\,\unit{\cubic\metre}\times2650 = \qty{2.65e15}{kg}$.

\textbf{20.} Het graniet is de smelt min zijn water: $2.65\times10^{15}/0.94 =
\qty{2.82e15}{kg}$ smelt.

\textbf{21.} Slechts $0.151\times0.70$ van de massa van de schist vertrekt als smelt:
$2.82\times10^{15}/(0.151\times0.70) = \qty{2.67e16}{kg}$.

\textbf{22.} $2.67\times10^{16}/2800 = \qty{9.5e12}{\cubic\metre}$, dat is
\qty{9.5e3}{\cubic\kilo\metre}, bijna tien keer de \omterm{def:b2:magma-emplacement:bodies}{pluton}.

\textbf{23.} $9.5\times10^3/5 = \qty{1.9e3}{\square\kilo\metre}$, een vierkant van ongeveer
\qty{44}{km} zijde.

\textbf{24.} Met $F = \qty{7.21}{\percent}$ is de bron $15.1/7.21 = 2.1$ keer
groter, ongeveer \qty{2.0e4}{\cubic\kilo\metre}. Het volume bronkorst dat werd gesmolten om
de \omterm{def:b2:magma-emplacement:bodies}{pluton} te maken, met de \omterm{def:b2:mantle-melting:fraction}{smeltfractie} van vraag~5, is
\textbf{\qty[text-series-to-math=true]{9.5e3}{\cubic\kilo\metre}}.
\end{solution}
